\section{Discussion and Validity Boundaries}
The experiments support a decision procedure, not a universal ranker: define the
exact storage population and task metric; test a small set of inexpensive
baselines; compare actual masks at relevant budgets under paired fields; and
spend additional computation only where measurable task benefit warrants it.
For the tested ResNet-8, per-example absolute gradients are useful allocation
proxies at several stress conditions. For KWS, corrected MSB-first is a strong
competitor, especially when a larger budget protects complete significant tiers.
For ToyCar, reconstruction sensitivity must be checked against file ordering;
good MSE prediction alone is inadequate evidence for AUC protection.

Several noncausal explanations remain possible. Gradient cancellation can
separate global and per-example scores, affine scales can separate numerical
jump from MSB ordering, and nonlinear inference can separate isolated damage
from multi-fault benefit. None of these facts identifies the dominant cause of
the cross-model results. The models differ in task, parameters, exposure count,
input distributions and surrogate quality. Weight sparsity is not a measured
causal explanation of redundancy, and the current study does not intervene on
architecture or retrain matched checkpoints.

Only one frozen checkpoint per architecture is evaluated. Fixed-cohort field
intervals do not describe variation over new models or new input populations.
The balanced candidate sample is small within each stratum. ToyCar has few
fault fields and an exposure ledger showing prior clean evaluation of all 248
official files and prior fault-outcome exposure of 200; the remaining 40 normal
and eight anomalous files do not provide a large balanced untouched cohort.
Further trials can improve Monte Carlo precision on reused inputs but cannot
create new input-level confirmation.

The static budget grid has only two interior points. A fair partial-tier
FKeras comparison remains the most consequential baseline gap. Surrogate
gradients are not integer-program derivatives, and the FKeras continuous
Hessian is a third, explicitly different numerical object. Exact isolated
measurements reduce ambiguity about finite effects but still do not certify
multi-fault allocation. Bounded conditional experiments address that distinction
locally; their uncertain and small gains against isolated allocation do not
justify claiming an effective general optimizer.

Finally, logical protected-bit count is deliberately the sole protection cost.
No area, power, energy, latency, synthesis or radiation-calibrated reliability
estimate follows from these software results. Runtime measurements concern the
analysis itself. Excluded biases, metadata and activations remain outside the
fault model rather than being declared universally safe. This restricted scope
makes the comparison auditable, but limits its transfer to other fault domains.
